% PDF-derived LaTeX transcription. Original manuscript source was unavailable.
% Published PDF SHA-256: 101d7dc4f8aaca7c3c14d0e360d6a6e3bab53b883cccee61eef4fad47ffb8187
% Extracted text SHA-256: df7c10333d58db18444ca7c23609a6c8bfd177ecf8619d34bffac74f3f5c502d
\documentclass[11pt]{article}
\usepackage[margin=1in]{geometry}
\usepackage{amsmath,amssymb}
\setlength{\emergencystretch}{3em}
\title{Stored-Byte Distinctness and Physical Population Accounting in Scientific Latent Datasets}
\author{}
\date{}
\begin{document}
\maketitle
\noindent\textbf{PDF-derived LaTeX transcription.} This editable text follows the PDF extraction. Refer to the published PDF for figures, table layouts, and mathematical notation.

% Extracted chunk 1

Stored-Byte Distinctness and Physical Population Accounting in
Scientific Latent Datasets

\section*{Abstract}

Stored-vector distinctness does not establish physical-case diversity or learned reconstruc-
tion. We examine retained chemistry, particle, density, and soft-body artifacts using numeric
storage readback, provenance joins, checkpoint metadata, and source-bound arithmetic. Nine
chemistry shards contain 515 finite, distinct vectors, while their manifest labels 494 examples
as failures. All 8400 particle rows and 1200 reference arrays have verified content identities and
joins; including the reference lane changes the measured byte accounting. A density checkpoint
contains no trained weights, and its layout expands scalar slots while reducing storage preci-
sion. A modular sampling driver yields 12600 parameter tuples rather than 3704400 nominal
combinations. Finite mathematical controls additionally show why cell-pooling or direct quan-
tized readback errors are not universal reconstruction lower bounds. These results separate byte
identity, population coverage, information accounting, and physical validity without claiming
original model or simulation replay.

\section*{1 Introduction}

A scientific training corpus has several sizes. Its number of files, stored vectors, molecular configu-
rations, trajectories, and independently varied physical cases need not agree. Likewise, a compact
stored representation can contain learned features, explicit coordinates, reference geometry, or meta-
data. A numerical total is scientifically interpretable only after its counting unit and information
boundary are specified.
This study examines retained artifacts from a chemistry cache, a particle-dynamics representa-
tion, a density representation, and a soft-body sampling driver. It contributes a concrete separation
of stored-byte distinctness, source identity, parameter-space coverage, and representation account-
ing. The principal evidence consists of reopened numeric storage, complete retained manifests, a
small checkpoint whose contents can be inspected without executing serialized objects, and source
expressions checked without running the supplied applications. The analysis does not retrain mod-
els, rerun physical trajectories, or establish chemical accuracy.
The cases reveal different issues. A deduplicated chemistry cache is byte-distinct but contains
predominantly failure-labelled examples. A particle vault has distinct dynamic and reference ten-
sors, yet its physical population cannot be reconstructed from the clean manifest alone. A density
checkpoint contains no learned weights, and its metadata-bearing layout expands the number of
scalar slots while reducing storage precision. A sampling driver’s correlated modular parameters
generate a much smaller catalogue than a product of nominal knob counts suggests. These findings
motivate population and information accounting as prerequisites for interpreting scientific recon-
struction results.
1


\section*{2 Related work}

Lee et al. distinguish exact substring duplication from approximate document duplication in
language-model datasets and describe methods for detecting both [ 1]. Their work shows why a
deduplication procedure must define the object and equivalence relation being compared. Here, the
objects are stored numeric rows and scientific source identities. We do not transfer their language-
model performance results to these physical datasets or treat a tensor hash as a chemical identity
test.
Kapoor and Narayanan’s July 2022 preprint distinguishes a lack of train–test separation from
a mismatch between an evaluation distribution and the distribution of scientific interest [ 2]. Its
discussion of nonindependent observations makes the intended scientific claim essential: dependence
can be appropriate for one target population and inappropriate for another. Our retained-data
analysis supplies case-specific accounting rather than a universal leakage verdict. Neither cited
work validates the private simulations, checkpoints, or physical claims examined here. We cite the
inspected preprint version explicitly; it is not the later journal version.

\section*{3 Materials and checking methods}

This is a retrospective analysis of retained artifacts; no additional model training, neural inference,
or physical simulations were performed. Supplementary record S1 contains the original retained
files, source snapshots, provenance receipts, and exact local checking inputs and results. Every
captured file is bound to its byte length and SHA-256 digest. Local checking calculations were
run in isolation; supplied application programs were inspected as source text, never imported or
executed. This separates measurements of retained artifacts from claims about how the original
model or simulator behaved.

\section*{3.1 Distinctness and provenance}

For stored rows bi and a deterministic content hash h, the measured count is
Ub = |\{h(bi) : i = 1, . . . , n\}| . (1)
Different hashes establish different byte inputs. Equal scientific meaning, independence of physical
cases, and validity of labels are separate questions. A source-identity map \ensuremath{\pi} may instead identify
a molecular system, simulation parameter case, trajectory, or observation time. The appropriate
population count is then
U\ensuremath{\pi} = |\{\ensuremath{\pi}(i) : i = 1, . . . , n\}| , (2)
provided the corresponding identity fields and their definitions are available. We do not infer U\ensuremath{\pi}
from Ub when those fields or the raw inputs are missing.
The chemistry data contain nine numeric shards. Their ZIP structures and a narrowly specified
tensor-storage schema were checked without pickle deserialization or execution of object construc-
tors. The underlying little-endian IEEE float16 storage was read directly. Each row was also
expanded to float32 to test the precise digest convention used by its manifest. For the particle
vault, all dynamic and reference NPY files were read as numeric uint16 containers holding BF16
bit patterns. Object arrays were excluded. Every dynamic-row BLAKE2b-128 digest was recom-
puted and compared with both the clean manifest and its separate provenance row.
2


\section*{3.2 Storage and checkpoint accounting}

Byte accounting uses actual file lengths and includes array headers. Dynamic and reference stor-
age are reported separately; reference files are amortized over the dynamic rows that use them.
An amortized accounting contribution is not a literal row-file size. Model weights, query coordi-
nates, normalization, and other required inputs must also be counted before claiming an end-to-end
compression ratio.
The density checkpoint was decoded through a primitive-only metadata grammar that permits
literal dictionaries, lists, strings, numbers, and references between these values. It permits no global
resolution, callable execution, persistent tensor identifier, or general object deserialization. Its
implementation digest was compared with the captured implementation source. Static inspection
separately checked whether the audit program constructs and saves an empty state dictionary.

\section*{3.3 Parameter catalogue and mathematical controls}

For the soft-body driver, the body of the parameter-selection function was structurally compared
with a fixed transcription of its arithmetic expressions. A separately written implementation of
those expressions enumerated 200000 sample indices and tested repetition over one complete period.
No supplied driver or simulator was executed. This establishes a parameter-tuple catalogue under
that formula, rather than a count of independent physical realizations.
Finally, finite mathematical examples test two proposed implications: that a cell-constant recon-
struction error bounds a decoder with position queries, and that raw BF16 readback error bounds
any learned reconstruction. These are counterexamples to universal statements, not experimental
measurements of a scientific model.

\section*{4 Results}


\section*{4.1 Chemistry: distinct stored rows and failure labels}

The retained deduplication summary reports 300096 upstream rows, 299581 removals, and 515 re-
tained rows. The partition arithmetic is consistent. Our direct measurements concern the retained
cache: nine shards contain 515 finite float16 vectors of length 2048, and all 515 stored rows have
distinct content hashes. All manifest digests match hashes of the stored values after float32 expan-
sion; hashes of the unexpanded float16 bytes are a different convention and do not match those
recorded digests.
The clean manifest also contains 515 distinct text strings. This is exact string distinctness,
without molecular canonicalization or a test of chemical equivalence. Its domain labels partition the
rows into 502 e-fuel, 10 battery, and 3 pharmaceutical examples. Its outcome categories partition
them into 21 success-labelled and 494 failure-labelled rows. The recorded audit flag is true for
every row. Thus, the flag and the outcome category describe different properties: a structurally
admissible retained record can represent a failed candidate. Neither the flag nor a unique vector
establishes a successful chemical result.
The upstream removal count remains a saved-summary result: we did not independently reopen
all 300096 original rows and their tensors. The retained-cache measurements therefore corroborate
its final contents and digest convention, not the entire historical deduplication execution.
3


\section*{4.2 Particle vault: distinct bytes, reference joins, and byte totals}

All 17 dynamic shards were reopened. Their 8400 BF16 rows have shape 64 \ensuremath{\times} 32, and all 8400
independently computed row digests match both manifest and provenance. All 1200 reference
arrays have shape 64 \ensuremath{\times} 16; their content digests match their reference-group identifiers. Each group
joins to seven dynamic rows with frame indices 0 through 6. Both dynamic rows and reference
contents are byte-distinct within their respective retained populations.
The distinction between records and physical populations remains essential. The clean manifest
has no populated sample identifier or raw path. Its separate provenance supplies 1200 reference
groups but only 480 distinct sample basenames across the indexed source roots. Basenames alone
cannot recover the complete source identity. The source audit describes a shared 48-cell generation
grid with repeated particle-subset exports, but the four original raw-data paths in the vault index
are no longer present at those recorded locations. We therefore retain that historical source-
audit interpretation separately from our direct byte and join measurements. We do not claim an
independently replayed count of 48 physical simulations or certify unseen-parameter generalization.
Table 1: Measured particle-vault storage, including NPY headers. Reference storage is amortized
over 8400 dynamic rows. These totals exclude model weights and any other required query or
normalization inputs.
Stored component Total bytes Bytes per dynamic row
Dynamic arrays 34408576 4096.259048
Reference arrays 2611200 310.857143
Combined 37019776 4407.116190
The vault’s recorded mean of 3994.819048 bytes per sample is reproduced by counting only the
16 full dynamic shards and dividing by all 8400 rows. It omits the final 852096-byte partial shard.
Including that shard gives the dynamic mean in the table. Adding the reference lane increases
the mean again. Distinct storage, a correct token count, and complete information accounting are
therefore three different checks; a pass on one does not imply a pass on the others.

\section*{4.3 Density representation: an empty checkpoint and carried constraints}

The actual 1983-byte checkpoint matches the digest recorded by its historical metadata. It con-
tains an empty state dictionary and no tensor-storage members. Its embedded implementation
digest matches the captured implementation source. The audit source constructs the empty state
dictionary explicitly before saving it. This supports a narrow conclusion about this artifact: the
checkpoint is metadata, not a collection of trained model parameters. It does not imply that every
other model in the programme is weightless.
The input density grid contains 483 = 110592 scalar values. Its layout stores 216 density patches
of 512 values and one additional 512-slot metadata token:
nstored = 217 \ensuremath{\times} 512 = 111104 , nstored \ensuremath{-} ninput
ninput
= 512
110592 \ensuremath{\simeq} 0.00462963. (3)
The scalar-slot count expands by approximately 0.463\%. FP32 input values require 442368 payload
bytes, whereas 111104 BF16 slots require 222208 payload bytes. This precision-based byte saving
must not be described as learned scalar compression. It also does not establish a physical-density
accuracy result.
4

The implementation reassembles the patch values, clips negative density, and rescales the result
using the carried electron count. Writing the clipped density as \ensuremath{\rho}+
i and voxel volume as v, the
relevant operation has the form
b\ensuremath{\rho}i = \ensuremath{\rho}+
i
Ne
max
[U+0010]
v P
j \ensuremath{\rho}+
j , \ensuremath{\varepsilon}
[U+0011] . (4)
Here, Ne is read from decoded metadata, and v is reconstructed by adding its two voxel-volume
metadata slots. The floor \ensuremath{\varepsilon} is the smallest positive normal value for the decoded density’s floating-
point dtype, implemented as torch.finfo(density.dtype).tiny. The supplied audit expands the
BF16 tokens to FP32 before decoding. For a finite positive integral above this floor, normalization
makes the resulting integral equal to the decoded Ne in exact arithmetic using the reconstructed v;
floating-point equality is not guaranteed, and decoded metadata need not equal its pre-quantization
values exactly. An identically zero clipped density remains zero for finite metadata. This normal-
ization can check implementation consistency, but it cannot independently establish spatial-density
accuracy or agreement with the original pre-storage electron count. We did not execute the original
framework or remeasure its density predictions.

\section*{4.4 Soft-body driver: correlated knobs and a finite catalogue}

The inspected driver selects a pair of geometry resolutions with period 25 and other parameters
through modular expressions with periods 9, 7, 6, and 8. Multiple nominal knobs share the same
underlying residue modulo 7. The joint catalogue therefore has period
L = lcm(25, 9, 7, 6, 8) = 12600 . (5)
Enumeration over 200000 consecutive indices yields exactly 12600 distinct full parameter tuples
and 25 geometry pairs; all tested tuples repeat after L. The naive product of the individual knob
counts, 3704400, overcounts this catalogue by a factor of 294.
This result concerns the specified parameter-selection formula. It does not prove that every tuple
produces an admissible or physically distinct trajectory, that different GPU executions are bitwise
identical, or that a different sampling driver has the same ceiling. The current captured formula
matches the transcription in the historical capacity note, but a complete historical execution and
environment binding was not replayed. Static source inspection also shows that the displayed time
integrator accepts no constitutive-stress argument; that signature fact is not itself a whole-program
momentum-balance measurement.

\section*{5 Interpretation and limitations}

These cases support a hierarchy of questions rather than a single corpus-quality score. Are the
original bytes readable and finite? Do the stored hashes and provenance join? What equivalence
relation defines a source case? What distribution does a held-out set represent? What information
reaches the decoder? Which quantities are predicted, and which are copied or enforced? The
answers must accompany the numerical result they qualify.

\section*{5.1 A weaker reconstruction family is not a universal bound}

Consider two distinct coordinates x1 and x2 assigned to the same spatial cell, with targets yi = xi.
A cell-constant reconstruction has positive squared error. A decoder that receives the individual
5

query xi can instead return xi exactly. Similarly, a quadratic target can be represented by a
nonlinear query decoder even when a per-cell affine fit has nonzero residual. Consequently, the
error of either weaker reconstruction family is not a universal lower bound for a decoder with
richer query information. The retained particle-study note explicitly withdraws its earlier universal
cell-pooling bound for this reason. We do not interpret the note’s obsolete bound language as a
validated theorem.
Raw quantized readback provides another useful baseline, but not a universal lower bound on
reconstruction error. For a scalar observation Q(X) and an available prior, the squared-error
estimator
g\ensuremath{\ast}(Q(X)) = E[X | Q(X)] (6)
can outperform returning the numerical center represented by Q(X). In our finite control, two
float32 values near 1.0001 and 1.0002 round to the same BF16 value 1. Their conditional mean
has smaller mean squared error than direct readback. This does not invert the quantizer for
arbitrary inputs; it demonstrates the role of the target distribution and available prior. It supplies
no empirical performance claim for any learned codec in the present collection.

\section*{5.2 Coverage and uncertainty}

The direct measurements cover retained storage, manifests, checkpoint metadata, and explicit
source expressions. They do not reproduce original neural inference or physical dynamics. Hash
verification authenticates the captured contents, not the truth of their labels. The chemistry
outcomes are recorded categories; no independent chemical oracle was run. The particle raw-input
paths are unresolved at their recorded locations, and distinct tensor rows cannot reconstruct the
missing physical-case map. Parameter enumeration is not a physical uncertainty estimate.
The four cases were selected from a broader research inventory. Unrecovered cross-vault overlap
and other specialist audit reports are not used as principal empirical results here. The source notes
contain later corrections and should not be flattened into an account in which every original claim
survived. Our two inspected methodological references provide context rather than a comprehensive
novelty search. No claim of global originality, full-corpus validation, or production readiness follows
from these results.

\section*{6 Conclusion}

A large count of distinct stored vectors can coexist with a much smaller parameter catalogue, carried
reference information, failure-labelled examples, or a checkpoint without trained weights. The
retained artifacts make these distinctions measurable: 515 chemistry vectors survive the reported
upstream deduplication; 8400 particle rows and 1200 reference arrays have verified byte identities
and joins; a density layout expands its scalar slots while using lower-precision storage; and a
modular driver contains 12600 parameter tuples rather than the naive 3704400 combinations.
The scientific interpretation depends on what each measurement establishes. Stored-byte dis-
tinctness is valuable, but it is not a substitute for physical-case identity or generalization evidence.
Explicit metadata and reference lanes can be useful, but they belong inside the information bound-
ary. Reconstruction baselines can diagnose a representation, but they do not become universal
lower bounds merely because they remove the learned model. Reporting these boundaries makes
6

the numerical results assessable without converting unresolved checks into either blanket rejection
or unsupported success.

\section*{7 Data and code availability}

Supplementary record S1 is a private evidence package containing the retained chemistry numeric
shards, particle dynamic and reference arrays, source-code snapshots, manifests, checkpoint meta-
data, and checking scripts, inputs, and results. Its 105 final aggregate and consistency checks
include measurements of the stored numeric rows and reference arrays.
No public archive, external-reviewer access mechanism, or redistribution permission is established
for the private project material or bundled literature. The captured data do not include a replay of
the original models, all upstream chemistry tensors, or the missing particle raw-input roots. The
full provenance and specific unresolved checks are recorded separately.

\section*{References}

[1] K. Lee, D. Ippolito, A. Nystrom, C. Zhang, D. Eck, C. Callison-Burch, and N. Carlini. Dedu-
plicating Training Data Makes Language Models Better. Proceedings of ACL, pp. 8424–8445
(2022). https://aclanthology.org/2022.acl-long.577/.
[2] S. Kapoor and A. Narayanan. Leakage and the Reproducibility Crisis in ML-based Science.
arXiv:2207.07048v1, July 14, 2022. https://arxiv.org/abs/2207.07048v1.
7

\end{document}
